%=====================================================================
\chapter{Themes, Layout and the Customer's Experience}\label{ch:ux}
%=====================================================================
\locator{2}{Part II --- Standing the Shop Up}

\begin{outcomes}
By the end of this chapter you will be able to:
\begin{tight}
  \item \textbf{Describe} how a theme is built from templates, and
        \textbf{override} one without editing the original.
  \item \textbf{Explain} what a layout is, and \textbf{place} a module on a
        chosen page without touching a template at all.
  \item \textbf{Measure} a page's weight and \textbf{say} which part of it is
        the weight.
  \item \textbf{Distinguish} the optimisations that are worth doing from the
        ones that are merely popular, using evidence.
  \item \textbf{Judge} a storefront change by its effect on the customer rather
        than by taste.
\end{tight}
\end{outcomes}

\begin{prereq}
\chref{catalogue} and \chref{variations}: a catalogue worth looking at.
Appendix~B's Twig section --- this chapter edits templates. HTML and CSS from
\emph{Web Technologies}.
\end{prereq}

\begin{keyterms}
theme $\bullet$ template $\bullet$ Twig $\bullet$ template override $\bullet$
layout $\bullet$ module $\bullet$ position $\bullet$ page weight $\bullet$
request count $\bullet$ render-blocking $\bullet$ minification $\bullet$
responsive design $\bullet$ above the fold $\bullet$ conversion
\end{keyterms}

\section{What a Theme Actually Is}

``Change the theme'' sounds like changing a colour. It is not. A theme in
OpenCart is a directory of Twig templates, and each one is responsible for one
piece of the page.

\begin{lstlisting}[language=bash]
docker compose exec shop \
  ls /var/www/html/catalog/view/template/product/
# category.twig  product.twig  search.twig  special.twig ...
\end{lstlisting}

The default theme lives in
\texttt{catalog/view/template/}. The important rule about it is one you should
adopt before you change a single line:

\begin{alertbox}[Never edit the default theme]
Edit \texttt{catalog/view/template/product/category.twig} directly and three
things happen. Your change is lost at the next upgrade. Nobody can tell what
you changed, because there is no original to compare against. And a second
developer, or you in March, will assume the file is stock and read it as
documentation of how OpenCart works.

\smallskip
Make a \term{child theme} instead: a directory of your own that contains
\emph{only} the templates you changed. OpenCart looks there first and falls
back to the default for everything else, so your directory \emph{is} the list
of your changes. A child theme with four files in it is a complete and honest
statement that you changed four things.
\end{alertbox}

\section{Layouts, or How to Change a Page Without Editing It}

A surprising amount of what people reach for a template to do can be done
without touching one.

\begin{definitionbox}[Layout, Module, Position]
A \term{module} is a self-contained block --- a category list, a banner, a
carousel, ``featured products''.

\smallskip
A \term{layout} says which modules appear on which \emph{kind} of page, and
where. \adminpath{Design, Layouts} lists them: Home, Category, Product,
Checkout, and so on.

\smallskip
A \term{position} is a slot in the page: \uiel{Column Left},
\uiel{Column Right}, \uiel{Content Top}, \uiel{Content Bottom}. A module is
placed in a position of a layout, with a sort order.

\smallskip
So ``put a banner above the products on every category page'' is three clicks
in \adminpath{Design, Layouts, Category} and no code at all.
\end{definitionbox}

\begin{conceptbox}[Ask this before editing a template]
\textbf{``Is this a layout change or a template change?''}

\smallskip
If you want something \emph{moved, added or removed} and it already exists as a
module, it is a layout change --- do it in the admin panel, where it is
visible to whoever comes next and survives upgrades.

\smallskip
Edit a template only when the \emph{markup itself} must differ: a field shown
that the theme does not show, a different arrangement within a block, a tag the
theme gets wrong. Then override exactly that file in your child theme and
nothing else.

\smallskip
The commonest waste of effort in this chapter's laboratory is a student editing
\texttt{category.twig} to do something \adminpath{Design, Layouts} would have
done.
\end{conceptbox}

\section{What a Page Weighs}

Here the chapter stops being about arrangement. A customer's experience of your
shop is mostly its speed, and speed is mostly \emph{weight} --- how many bytes
and how many separate requests the browser must complete before the page is
usable.

Note that this is a different question from \chref{catalogue}'s. That chapter
measured how long the \emph{server} took to produce the HTML. This one measures
everything that happens afterwards, and it is usually larger.

\begin{measured}[What the Storefront Weighs]
Produced by \texttt{code/m08\_pageweight.py} on the machine in Appendix~A. The
harness fetches a page, finds every stylesheet, script and image the HTML
names, requests each one once, and totals them.

\rowcolors{2}{white}{lightbg}
\begin{longtable}{@{}L{2.0cm}L{1.9cm}L{1.6cm}L{1.6cm}L{1.6cm}L{1.6cm}L{2.0cm}@{}}
\toprule
\rowcolor{primary!15}
\textbf{Page} & \textbf{Requests} & \textbf{HTML} & \textbf{CSS} &
\textbf{JS} & \textbf{Images} & \textbf{Total} \\
\midrule
\endhead
Home & 25 & 32\,KB & 348\,KB & 223\,KB & 293\,KB & \textbf{895\,KB} \\
Category & 20 & 46\,KB & 348\,KB & 223\,KB & 122\,KB & \textbf{738\,KB} \\
Product & 18 & 54\,KB & 354\,KB & 243\,KB & 61\,KB & \textbf{713\,KB} \\
\bottomrule
\end{longtable}

\smallskip
\textbf{The framework is the page.} CSS and JavaScript come to about
\textbf{570\,KB on every page}, and they do not vary with the content ---
the same bytes on the home page, a category of 50{,}000 products and a single
product. On the product page they are \textbf{84\% of the weight}. The HTML
that the whole of \chref{catalogue} was spent optimising is 6--7\%.

\smallskip
\textbf{Images are not the problem here, although everyone optimises them
first.} They are 33\% of the home page, 16\% of a category page and 9\% of a
product page. The largest single image in the demonstration catalogue is
13.9\,KB. You could halve every image on the category page and save 61\,KB ---
about 8\%.

\smallskip
\textbf{The single largest file is 265\,KB of unminified CSS.}
\texttt{bootstrap.css} ships at 271{,}637 bytes, and \textbf{no minified
version ships with it} --- there is no \texttt{bootstrap.min.css} anywhere in
the tree. A minified build of the same framework is around a tenth of that.
That one file is 36\% of the category page.

\smallskip
\textbf{This measurement undercounts, and you should know how.} It follows only
assets the \emph{HTML} names. Fonts are requested by the CSS, so they are
invisible to it: the FontAwesome directory alone holds 248\,KB of
\texttt{.woff2} files, of which a browser fetches whichever subsets it needs.
The totals above are therefore a floor, not a figure. A tool that drove a real
browser would see more; this one is honest about seeing less.
\end{measured}

\begin{conceptbox}[The order to optimise in, which is not the popular order]
From the numbers above, for this shop:

\begin{tightnum}
  \item \textbf{Minify the CSS.} One file, 265\,KB to roughly 27\,KB. The
        largest single saving available, and it changes nothing a customer can
        see.
  \item \textbf{Drop what you do not use.} Bootstrap and FontAwesome are
        general-purpose libraries; a shop uses a fraction of either. Removing
        the unused parts is the second-largest saving and the fiddliest.
  \item \textbf{Cut the request count.} Twenty requests is twenty round trips.
        On a fast link this hardly matters; on a phone on a slow connection in
        Bahawalpur it is most of the wait.
  \item \textbf{Then images} --- resize to the size actually displayed, and
        serve modern formats.
\end{tightnum}

\smallskip
Most advice reverses this list, because images are the thing people can see and
a 2\,MB photograph is a vivid story. On \emph{this} shop images are 16\% and
the framework is 77\%. On a shop with photographic product shots the balance
moves. That is the point: the order depends on the measurement, and the
measurement takes two minutes.
\end{conceptbox}

\begin{alertbox}[Render-blocking, and why position matters as much as size]
A stylesheet in the \texttt{<head>} is \term{render-blocking}: the browser will
not paint \emph{anything} until it has fetched and parsed it. So 348\,KB of CSS
is not merely 348\,KB of transfer --- it is 348\,KB standing between the
customer and the first pixel.

\smallskip
Scripts block too, unless marked \texttt{defer} or \texttt{async}. A
\texttt{<script src>} in the head stops the parser dead while it downloads.

\smallskip
This is why page weight and page time are not the same measurement, and why
``it is only 200\,KB'' can still be the reason a shop feels slow.
\end{alertbox}

\section{Judging a Change by the Customer}

The hardest discipline in this chapter is not technical. It is resisting the
urge to change the storefront because you would prefer it differently.

\begin{examplebox}[Three changes, and how you would know if they helped]
\textbf{``Make the Add to Cart button green.''} You cannot know without an
experiment: show half your visitors one colour and half the other, and compare
the proportion who buy. With thirty visitors a day you will wait months for an
answer, which is itself the answer --- a shop this size should spend its
effort elsewhere.

\smallskip
\textbf{``Show the delivery charge on the product page.''} This you can reason
about. \chref{what}'s laboratory found that shops which hide the delivery
charge until checkout lose customers there; it is the most-cited reason for
abandonment. No experiment needed.

\smallskip
\textbf{``Minify the CSS.''} Measurable directly, in bytes, with no customers
involved at all. 265\,KB to 27\,KB is a fact, and it cannot make anything
worse.

\smallskip
\textbf{The pattern.} Prefer changes you can verify without an experiment you
cannot afford to run. Most small shops cannot generate enough traffic to
measure a colour, and should spend their effort on the things that are
measurable in bytes or known from everybody else's evidence.
\end{examplebox}

\begin{pitfall}
\begin{tight}
  \item \textbf{Editing the default theme.} Lost at upgrade, invisible to the
        next person. Override in a child theme.
  \item \textbf{Editing a template to do a layout's job.} If the block exists
        as a module, move it in \adminpath{Design, Layouts}.
  \item \textbf{Optimising images first because they are visible.} On this shop
        that addresses 16\% of the weight and ignores 77\%.
  \item \textbf{Assuming a framework file is minified.} \texttt{bootstrap.css}
        here is 265\,KB and there is no minified copy in the tree.
  \item \textbf{Measuring page weight by counting every \texttt{<img>} tag.}
        The same thumbnail on twenty cards is one request; counting it twenty
        times overstates the weight several-fold.
  \item \textbf{Forgetting fonts.} They are requested by CSS, so they are
        invisible to any harness that reads only the HTML --- including this
        chapter's.
  \item \textbf{Redesigning on taste with no way to tell whether it helped.}
\end{tight}
\end{pitfall}

\begin{lab}[Weigh your shop, then make it lighter]
\begin{lstlisting}[language=bash]
# --- 1. Weigh three pages -------------------------------------------
cd code
python m08_pageweight.py --compare

# --- 2. Find the biggest single file --------------------------------
# Predict it before you look. Most people guess an image.
docker compose exec shop sh -c \
  'ls -laS /var/www/html/catalog/view/stylesheet/ | head -5'

# --- 3. Make a child theme ------------------------------------------
# Admin: Design > Theme Editor. Override ONE template -- the category
# product card -- and add the model number under each product name.
# Your override directory is now the list of your changes.

# --- 4. Do a layout change with no code -----------------------------
# Admin: Design > Layouts > Category. Add a module to Content Top.
# Reload a category page. No template was edited.

# --- 5. Weigh it again ----------------------------------------------
# Did your two changes move the number? By how much, and which one?
python m08_pageweight.py

# --- 6. Make a real saving ------------------------------------------
# Minify bootstrap.css and serve the minified copy instead. Any
# minifier will do. Record the before and after byte counts.

# --- 7. Weigh it a third time, and account for the difference -------
# The saving should be most of 240 KB. If it is not, find out what
# else changed -- do not just record the number.
python m08_pageweight.py

# --- 8. Say what you would do next, and why -------------------------
# One paragraph, ordered by measured saving rather than by what is
# easiest to see. Chapter 20 returns to this with caching.
\end{lstlisting}

\smallskip
\textbf{What to notice.} Step 2 against step 1: almost everyone predicts an
image and the answer is a stylesheet. That gap between where people look for
weight and where weight is, is the whole chapter.
\end{lab}

\begin{checkpoint}
\begin{tightnum}
  \item A product page is 713\,KB, of which 597\,KB is CSS and JavaScript. A
        colleague proposes compressing all the product photographs. What is the
        best case for that change, and what would you propose instead?
  \item Distinguish a layout change from a template change, and give one
        example of each for the request ``show a trust badge under the Add to
        Cart button on every product page''.
  \item This chapter's harness reports 738\,KB for a category page, and a
        browser's developer tools report more. Give two reasons the harness is
        lower, and say which of them it could fix.
\end{tightnum}
\end{checkpoint}

\begin{chaptersummary}
\begin{tight}
  \item A theme is a directory of Twig templates. Never edit the default:
        override only the files you change in a child theme, so that the
        directory is the list of your changes and survives upgrades.
  \item A layout places modules in positions on a kind of page. Moving, adding
        or removing an existing block is a layout change made in the admin
        panel, not a template edit.
  \item Measured: the storefront weighs 713--895\,KB per page across 18--25
        requests. CSS and JavaScript are about 570\,KB of that \emph{on every
        page} and do not vary with content --- 84\% of the product page.
  \item Images are 9--33\% depending on the page, although they are what people
        optimise first.
  \item The largest single file is \texttt{bootstrap.css} at 265\,KB,
        unminified, with no minified copy anywhere in the tree.
  \item The harness undercounts: it follows only assets named in the HTML, so
        fonts requested by CSS (248\,KB in one directory) are invisible to it.
        Know what your instrument cannot see.
  \item Weight and time are not the same: a stylesheet in the head is
        render-blocking and stands between the customer and the first pixel.
  \item Prefer changes you can verify without an experiment you cannot afford
        to run. A small shop cannot measure a button colour and can measure
        240\,KB.
\end{tight}
\end{chaptersummary}

\begin{reviewq}
  \item \textbf{(MCQ)} On the measured product page, CSS and JavaScript are
        what share of the weight? (a)~9\%; (b)~about a third; (c)~84\%;
        (d)~it depends on the catalogue size.
  \item \textbf{(MCQ)} ``Show the manufacturer's name under the product title''
        where the theme does not show it at all is: (a)~a layout change;
        (b)~a template override; (c)~a module setting; (d)~a database change.
  \item \textbf{(Short)} Why is a child theme better than editing the default,
        giving two reasons that are not about upgrades.
  \item \textbf{(Short)} Explain render-blocking, and why page weight and page
        time are different measurements.
  \item \textbf{(Short)} This chapter's harness counts each unique asset once
        rather than once per mention. Explain what would go wrong otherwise,
        with a number.
  \item \textbf{(Applied)} You are given a shop whose category page is 2.1\,MB:
        HTML 40\,KB, CSS 90\,KB, JS 150\,KB, images 1.82\,MB across 40
        requests. Give your optimisation order with a predicted saving for each
        step, and say how this differs from the order you would use on the shop
        in this chapter and why.
  \item \textbf{(Applied)} Design a measurement that would tell you whether
        showing the delivery charge on the product page increases or decreases
        completed orders. State what you would record, how long you would need
        to run it, and --- honestly --- whether a shop taking thirty orders a
        day could reach a conclusion at all.
\end{reviewq}
